\section{Architecture and Component Map}

\begin{figure}[htbp]
\centering
\resizebox{0.98\textwidth}{!}{%
\begin{tikzpicture}[x=1mm,y=1mm]
  % Operator and serving boundary.
  \node[client,minimum width=40mm] (browser) at (24,58)
    {React 19 browser UI\\search, inspection, workspace};
  \node[service,minimum width=50mm] (api) at (104,58)
    {FastAPI orchestration\\KIS / TRAKE / media / workspace};
  \node[artifact,minimum width=39mm] (state) at (165,58)
    {Canonical local state\\FrameStore / indexes / SQLite};

  % Canonical online execution path.
  \node[process,minimum width=31mm] (corpus) at (25,35)
    {Corpus and identity\\assets + evidence};
  \node[process,minimum width=33mm] (retrieval) at (65,35)
    {Multimodal retrieval\\FAISS + BM25 + fusion};
  \node[process,minimum width=35mm] (temporal) at (104,35)
    {Temporal alignment\\planner + monotonic DP};
  \node[process,minimum width=34mm] (materialize) at (144,35)
    {Task materialization\\KIS frame / TRAKE path};

  % Replaceable compute and artifact-publication boundary.
  \node[service,minimum width=48mm] (gpu) at (30,7)
    {Optional private inference\\SigLIP2 / BGE-M3 / Qwen / Florence-2};
  \node[service,minimum width=50mm] (offline) at (105,7)
    {Offline preparation worker\\FFmpeg + GPU enrichment / indexing};
  \node[artifact,minimum width=34mm] (s3) at (165,7)
    {S3 artifact transport\\versioned bundles};

  \draw[flowarrow] (browser) -- (api);
  \draw[flowarrow] (api) -- (state);
  \draw[flowarrow] (api.south) -- (temporal.north);
  \draw[flowarrow] (corpus) -- (retrieval);
  \draw[flowarrow] (retrieval) -- (temporal);
  \draw[flowarrow] (temporal) -- (materialize);

  \draw[dataarrow] ([xshift=-15mm]api.south) -- ++(0,-6mm)
    -| ([xshift=14mm]gpu.north);
  \draw[dataarrow] (offline) -- (gpu);
  \draw[dataarrow] (offline) -- (s3);
  \draw[dataarrow] (s3.north) -- (state.south);

  \begin{scope}[on background layer]
    \node[layer,fit=(corpus)(retrieval)(temporal)(materialize),
      label={[font=\footnotesize\bfseries,text=ReportSecondary]
      below:Canonical retrieval and temporal path}] {};
  \end{scope}
\end{tikzpicture}%
}
\caption{System architecture and trust boundary. Solid arrows are online control/data flow; dashed arrows are optional inference or artifact transport. Retrieval state is local and canonical; the GPU service is replaceable. Every box has a distinct label and border in addition to color.}
\label{fig:architecture}
\end{figure}

\begin{table}[htbp]
\centering
\caption{Implemented model and algorithm inventory. ``Active'' means present in the documented preparation or online path; optional services are never implied to be mandatory at serving time.}
\label{tab:models}
\small
\begin{tabularx}{\textwidth}{@{}>{\bfseries}p{0.15\textwidth}p{0.26\textwidth}Xp{0.12\textwidth}@{}}
\toprule
Capability & Model / algorithm & Produced evidence and downstream use & Status \\
\midrule
Caption & Qwen/Qwen3-VL-8B-Instruct & Frame caption; first \FrameContext{} section and inspectable specialist artifact & Offline active \\
Visible text & Florence-2-base-ft & OCR regions, frame summary, quality; gated \FrameContext{} section and BM25 OCR field & Offline active \\
Objects & YOLOE-26l-seg & Detections and normalized object summary; \FrameContext{} section and literal filter & Offline active \\
Speech & Qwen3-ASR-1.7B-hf & Timestamped transcript segments; BGE-M3 segment index and BM25 ASR field & Offline active \\
Speakers & pyannote community-1 & Diarization attached to transcript segments & Configurable \\
Visual encoder & SigLIP2 base patch16-224 & 768-D image vectors and visual-language query vectors & Active \\
Text encoder & BAAI/bge-m3 & 1024-D \FrameContext{} and ASR-segment vectors & Active \\
Dense search & FAISS \code{IndexFlatIP} & Exact inner product over L2-normalized vectors; cosine-equivalent ranking & Active \\
Lexical search & Fielded BM25 & Title/caption/OCR/ASR literal evidence & Active \\
Query preparation & Qwen/Qwen3-4B & Five optional candidate rewrites selected explicitly by the client & Optional \\
Temporal decoding & Prefix-max monotonic DP & Strict same-video event paths, alternatives, and canonical backpointers & KIS/TRAKE active \\
VLM reranking & Qwen3-VL reranker adapter & Experimental package capability; absent from public request contract and UI & Not wired \\
\bottomrule
\end{tabularx}
\end{table}

\clearpage
\section{Offline Preparation and Index Construction}

\begin{figure}[htbp]
\centering
\resizebox{0.98\textwidth}{!}{%
\begin{tikzpicture}[x=1mm,y=1mm]
  % A compact tree uses vertical space so every stage remains readable and
  % connectors never pass through another node.
  \node[service,minimum width=28mm] (raw) at (15,55) {Raw videos};
  \node[service,minimum width=28mm] (provided) at (15,38) {Organizer keyframes\\+ mapping};
  \node[process,minimum width=30mm] (extract) at (55,55)
    {C++/FFmpeg sampler\\nearest frame at $k\!\times\!1$ s};
  \node[process,minimum width=30mm] (import) at (55,38)
    {Order-preserving import\\mapped IDs and times};
  \node[artifact,minimum width=30mm] (frame) at (95,46.5)
    {Canonical FrameStore\\IDs, time, assets};

  \node[process,minimum width=26mm] (siglip) at (130,57) {SigLIP2\\image encoding};
  \node[artifact,minimum width=29mm] (vindex) at (165,57) {Visual FAISS\\768-D frames};
  \node[process,minimum width=30mm,minimum height=13mm] (specialists) at (130,40)
    {Qwen3-VL caption\\Florence-2 OCR\\YOLOE objects};
  \node[artifact,minimum width=28mm] (context) at (165,40)
    {FrameContext\\caption / text / objects};

  \node[artifact,minimum width=32mm] (bm25) at (120,23)
    {Fielded BM25\\title / caption / OCR / ASR};
  \node[process,minimum width=25mm] (contextbge) at (165,23)
    {BGE-M3\\text encoding};
  \node[artifact,minimum width=29mm] (cindex) at (165,6)
    {Context FAISS\\1024-D frames};

  \node[process,minimum width=30mm] (asr) at (50,6)
    {Qwen3-ASR + pyannote\\timestamped segments};
  \node[process,minimum width=25mm] (asrbge) at (85,6)
    {BGE-M3\\segment encoding};
  \node[artifact,minimum width=29mm] (aindex) at (120,6)
    {ASR FAISS\\1024-D segments};

  \draw[flowarrow] (raw) -- (extract);
  \draw[flowarrow] (extract) -- (frame);
  \draw[flowarrow] (provided) -- (import);
  \draw[flowarrow] (import) -- (frame);

  \draw[flowarrow] (frame) -- (siglip);
  \draw[flowarrow] (siglip) -- (vindex);
  \draw[flowarrow] (frame) -- (specialists);
  \draw[flowarrow] (specialists) -- (context);
  \draw[flowarrow] (context) -- (contextbge);
  \draw[flowarrow] (contextbge) -- (cindex);

  \draw[flowarrow] (raw.west) -- ++(-4mm,0) |- (asr.west);
  \draw[flowarrow] (asr) -- (asrbge);
  \draw[flowarrow] (asrbge) -- (aindex);

  \draw[dataarrow] (frame.south) -- ++(0,-12mm)
    -| node[pos=0.70,above,font=\scriptsize]{title metadata}
    ([xshift=-7mm]bm25.north);
  \draw[dataarrow] (context.south) -- ++(0,-4mm)
    -| ([xshift=7mm]bm25.north);
  \draw[dataarrow] (asr.north) -- ++(0,5mm) -| (bm25.south);
\end{tikzpicture}%
}
\caption{Two ingestion routes converge on one canonical frame identity. Solid arrows show model or encoding stages; dashed arrows feed the fielded lexical index. The frame-native visual/context indexes and segment-native ASR index preserve their native evidence geometry rather than forcing all sources into one lossy representation.}
\label{fig:preparation}
\end{figure}

\begin{table}[htbp]
\centering
\caption{Observed materialized artifact snapshot. Counts document system scale, not retrieval effectiveness. Missing provenance fields are reported rather than inferred.}
\label{tab:artifacts}
\small
\begin{tabularx}{\textwidth}{@{}>{\bfseries}p{0.20\textwidth}p{0.15\textwidth}p{0.20\textwidth}X@{}}
\toprule
Artifact & Entities & Representation & Observed contract / caveat \\
\midrule
Canonical FrameStore & 873 videos; 470{,}804 frames & Parquet + manifest; 1-s targets & \code{custom-raw1fps-v1}; no duplicate submission-coordinate groups \\
Caption / OCR / objects / context & 470{,}804 rows each & Specialist Parquet + deterministic \FrameContext{} & 2{,}689{,}621 OCR regions and 5{,}654{,}993 object detections \\
Visual index & 470{,}804 frames & 768-D float32, L2, flat IP & Model name recorded; serialized model revision/source/config fingerprints absent \\
Vietnamese context index & 470{,}804 frames & 1024-D float32, L2, flat IP & Dataset label differs from FrameStore; revision/retrieval-source/fingerprints absent \\
ASR segment index & 40{,}156 intervals & 1024-D float32, L2, flat IP & BGE-M3 revision recorded; half-open temporal geometry retained \\
BM25 index & 470{,}804 documents & Four sparse field matrices & $k_1=1.5$, $b=0.75$; current caption vocabulary size is zero \\
\bottomrule
\end{tabularx}
\end{table}

\begin{table}[H]
\centering
\caption{Key deterministic policies and defaults used in this report.}
\label{tab:parameters}
\small
\begin{tabularx}{\textwidth}{@{}>{\bfseries}p{0.18\textwidth}p{0.26\textwidth}X@{}}
\toprule
Stage & Parameter & Policy \\
\midrule
Sampling & $\Delta=1000$ ms & Select nearest decoded timestamp; ties choose preceding frame; targets lie in $[0,D)$ \\
FrameContext & budgets $80/80/40$; OCR $\geq0.5$ & Ordered caption, visible-text, object sections; preserve lineage and availability \\
Dense bundles & float32, L2, flat IP & Exact cosine-equivalent search; vectors, mappings, postings, and timestamps stored together \\
ASR projection & maximum gap 5000 ms & Prefer frames inside $[a,b)$; otherwise nearest midpoint within bound for detached search \\
Detached fusion & RRF constant 60 & Weighted union of modality top-$K$ lists; visual source required \\
Temporal fusion & adaptive, coverage-aware & Quantile calibration, reliability gating, cue routing, per-frame renormalization \\
Alignment & $\lambda=10^{-5}$; power 1 & Strict monotonic path; prefix-max recurrence; latest predecessor on ties \\
Alternatives & 5 paths/video; 30 s separation & Suppress near-duplicate final-event moments before global ranking \\
\bottomrule
\end{tabularx}
\end{table}

\clearpage
\section{Online Fusion and Temporal Task Orchestration}

\begin{figure}[htbp]
\centering
\resizebox{0.98\textwidth}{!}{%
\begin{tikzpicture}[node distance=6mm and 7mm,scale=0.95,transform shape]
  \node[client,minimum width=29mm] (query) {Query input\\plain KIS or E1..En};
  \node[process,minimum width=31mm,right=8mm of query] (plan) {Planner / optional\\selected rewrite};

  \node[process,minimum width=27mm,right=9mm of plan,yshift=19mm] (visual) {SigLIP2\\visual scores};
  \node[process,minimum width=27mm,right=9mm of plan,yshift=6mm] (context) {BGE-M3\\context scores};
  \node[process,minimum width=27mm,right=9mm of plan,yshift=-7mm] (asr) {BGE-M3\\ASR interval scores};
  \node[process,minimum width=27mm,right=9mm of plan,yshift=-20mm] (lexical) {BM25\\field scores};

  \node[service,minimum width=36mm,right=11mm of context,yshift=-6mm] (fusion) {Calibration + routing\\reliability $\times$ coverage};
  \node[artifact,minimum width=31mm,right=10mm of fusion] (matrix) {Per-video matrix\\$F_{e,t}$};
  \node[service,minimum width=35mm,right=10mm of matrix] (dp) {Prefix-max DP\\strict chronological paths};

  \node[client,minimum width=30mm,below=13mm of dp,xshift=-19mm] (kis) {KIS output\\representative frame + path};
  \node[client,minimum width=30mm,right=8mm of kis] (trake) {TRAKE output\\complete ordered path};

  \draw[flowarrow] (query) -- (plan);
  \draw[flowarrow] (plan) -- (visual);
  \draw[flowarrow] (plan) -- (context);
  \draw[flowarrow] (plan) -- (asr);
  \draw[flowarrow] (plan) -- (lexical);
  \draw[flowarrow] (visual) -- (fusion);
  \draw[flowarrow] (context) -- (fusion);
  \draw[flowarrow] (asr) -- (fusion);
  \draw[flowarrow] (lexical) -- (fusion);
  \draw[flowarrow] (fusion) -- (matrix);
  \draw[flowarrow] (matrix) -- (dp);
  \draw[flowarrow] (dp) -- (kis);
  \draw[flowarrow] (dp) -- (trake);

  \node[note,below=5mm of fusion,align=center] {Full-corpus score fusion (not top-$K$ RRF)};
\end{tikzpicture}%
}
\caption{Shared KIS/TRAKE online path. Original event text remains available to lexical scoring; any rewrite is explicit. Component support and confidence are applied before the same monotonic decoder produces task-specific materializations.}
\label{fig:temporal}
\end{figure}

\begin{figure}[htbp]
\centering
\begin{tikzpicture}[x=12mm,y=9mm]
  \foreach \x/\lab in {0/$t_0$,1/$t_1$,2/$t_2$,3/$t_3$,4/$t_4$,5/$t_5$,6/$t_6$,7/$t_7$}
    \node[font=\scriptsize,text=ReportSecondary] at (\x,-0.55) {\lab};
  \foreach \y/\event in {0/$E_1$,1/$E_2$,2/$E_3$}
    \node[font=\small\bfseries,anchor=east] at (-0.45,\y) {\event};
  \foreach \x in {0,...,7}
    \foreach \y in {0,...,2}
      \node[draw=ReportGrid,circle,minimum size=4.5mm,inner sep=0pt,fill=white] (n\y\x) at (\x,\y) {};
  \node[draw=ReportBlue,fill=ReportBlueTint,circle,minimum size=5.4mm,inner sep=0pt] (p00) at (1,0) {};
  \node[draw=ReportBlue,fill=ReportBlueTint,circle,minimum size=5.4mm,inner sep=0pt] (p11) at (4,1) {};
  \node[draw=ReportBlue,fill=ReportBlueTint,circle,minimum size=5.4mm,inner sep=0pt] (p22) at (6,2) {};
  \draw[flowarrow,draw=ReportBlue,line width=1.2pt] (p00) -- node[above,sloped,font=\scriptsize]{$u<t$} (p11);
  \draw[flowarrow,draw=ReportBlue,line width=1.2pt] (p11) -- node[above,sloped,font=\scriptsize]{$u<t$} (p22);
  \draw[decorate,decoration={brace,amplitude=3pt,mirror},draw=ReportOrange]
    (1,-0.92) -- node[below=3pt,font=\scriptsize,text=ReportSecondary]{gap penalty $\lambda(\tau_t-\tau_u)$} (4,-0.92);
  \node[note,anchor=west] at (8.0,1.55) {Selected nodes maximize\\event evidence plus the best\\valid prefix predecessor.};
  \node[note,anchor=west] at (8.0,0.25) {Backpointers recover one\\strict same-video path.};
\end{tikzpicture}
\caption{Dynamic-programming interpretation. Each event selects a strictly later canonical frame. Prefix maxima replace an all-pairs predecessor search, while backpointers preserve exact frame identity.}
\label{fig:dp-explainer}
\end{figure}

\clearpage
\section{Live Operator Interface}

The screenshots below were captured from the running React/FastAPI application at
1440$\times$900. To keep this capture isolated and reproducible on the report host,
the runtime used an identity-preserving subset of 21 videos and 3{,}807 frames from
the documented corpus, together with their original visual vectors and specialist
evidence; only visual dense retrieval was enabled in this capture profile. Because
the original source video was unavailable locally, inspector playback uses a
clearly scoped proxy assembled from the canonical 1-FPS frames of
\code{L30\_V060}; it exercises the real range-streaming and timeline code but is not
presented as the source video. Corpus totals and algorithms elsewhere in this report
come from the full artifact manifests, not from this UI fixture. Numbered markers,
component borders, and captions make the evidence readable without relying on color
alone.

\begin{figure}[htbp]
\centering
\begin{tikzpicture}
  \node[anchor=south west,inner sep=0] (image) at (0,0)
    {\screenshotplaceholder{figures/screenshots/kis-results.png}{0.96\textwidth}};
  \begin{scope}[x={(image.south east)},y={(image.north west)}]
    \node[callout] at (0.18,0.88) {1};
    \node[callout] at (0.09,0.67) {2};
    \node[callout] at (0.61,0.80) {3};
    \node[callout] at (0.31,0.38) {4};
  \end{scope}
\end{tikzpicture}
\caption{KIS retrieval workspace. \calloutlabel{1} Natural-language query and mode-sensitive input; \calloutlabel{2} Top-$K$, dense/BM25 controls, and shared answer files; \calloutlabel{3} latency and result summary; \calloutlabel{4} canonical result cards with an expanded event alignment and submission action.}
\label{fig:ui-kis}
\end{figure}

\begin{figure}[htbp]
\centering
\begin{tikzpicture}
  \node[anchor=south west,inner sep=0] (image) at (0,0)
    {\screenshotplaceholder{figures/screenshots/trake-results.png}{0.96\textwidth}};
  \begin{scope}[x={(image.south east)},y={(image.north west)}]
    \node[callout] at (0.24,0.87) {1};
    \node[callout] at (0.48,0.58) {2};
    \node[callout] at (0.28,0.77) {3};
  \end{scope}
\end{tikzpicture}
\caption{TRAKE ordered-event retrieval. \calloutlabel{1} Sequential \code{E1..En} specification; \calloutlabel{2} one same-video path displayed in event order; \calloutlabel{3} path-level submission action, which emits one ordered multi-frame CSV row rather than independent frame submissions.}
\label{fig:ui-trake}
\end{figure}

\begin{figure}[htbp]
\centering
\begin{tikzpicture}
  \node[anchor=south west,inner sep=0] (image) at (0,0)
    {\screenshotplaceholder{figures/screenshots/video-inspector.png}{0.96\textwidth}};
  \begin{scope}[x={(image.south east)},y={(image.north west)}]
    \node[callout] at (0.34,0.56) {1};
    \node[callout] at (0.30,0.10) {2};
    \node[callout] at (0.80,0.70) {3};
    \node[callout] at (0.82,0.40) {4};
  \end{scope}
\end{tikzpicture}
\caption{Video evidence inspector. \calloutlabel{1} media preview positioned at the retrieved moment; \calloutlabel{2} custom timeline and precise seek controls; \calloutlabel{3} live canonical/submission coordinates; \calloutlabel{4} available caption, OCR, object, and other inspectable evidence used for human verification.}
\label{fig:ui-player}
\end{figure}

\begin{figure}[htbp]
\centering
\begin{tikzpicture}
  \node[anchor=south west,inner sep=0] (image) at (0,0)
    {\screenshotplaceholder{figures/screenshots/collaborative-workspace.png}{0.96\textwidth}};
  \begin{scope}[x={(image.south east)},y={(image.north west)}]
    \node[callout] at (0.30,0.72) {1};
    \node[callout] at (0.82,0.78) {2};
    \node[callout] at (0.82,0.59) {3};
    \node[callout] at (0.82,0.48) {4};
  \end{scope}
\end{tikzpicture}
\caption{Collaborative Workspace. \calloutlabel{1} user-scoped query history and replay; \calloutlabel{2} manual video/timestamp opening; \calloutlabel{3} shared SQLite-backed submission files with distinct validated and filled states; \calloutlabel{4} multi-file CSV ZIP export.}
\label{fig:ui-workspace}
\end{figure}

\clearpage
\section{Traceability and Deployment Evidence}

\begin{table}[htbp]
\centering
\caption{Coverage matrix for every subsystem presented in the four-page body. UI behavior is evidenced by a live capture; nonvisual computation is evidenced by a diagram and/or manifest table.}
\label{tab:coverage}
\small
\begin{tabularx}{\textwidth}{@{}>{\bfseries}p{0.22\textwidth}Xp{0.30\textwidth}@{}}
\toprule
Presented subsystem & Appendix evidence & Primary implementation boundary \\
\midrule
End-to-end architecture / services & Fig.~\ref{fig:architecture}; Table~\ref{tab:models} & \pathcode{src/hcmai/orchestration}, \pathcode{src/hcmai/app.py}, \pathcode{llm/} \\
Keyframe ingestion / identity & Fig.~\ref{fig:preparation}; Table~\ref{tab:parameters} & \pathcode{offline/keyframes}, \pathcode{offline/ingestion} \\
Caption, OCR, objects, ASR, FrameContext & Fig.~\ref{fig:preparation}; Tables~\ref{tab:models}--\ref{tab:artifacts} & \pathcode{offline/enrichment/context} \\
FAISS / BM25 / fusion & Figs.~\ref{fig:preparation} and~\ref{fig:temporal}; Tables~\ref{tab:artifacts}--\ref{tab:parameters} & \pathcode{src/hcmai/retrieval} \\
Temporal DP / KIS / TRAKE & Figs.~\ref{fig:temporal} and~\ref{fig:dp-explainer} & \pathcode{src/hcmai/temporal}, \pathcode{orchestration/workflows} \\
KIS operator flow & Fig.~\ref{fig:ui-kis} & \pathcode{frontend/src/features/search} \\
TRAKE operator flow & Fig.~\ref{fig:ui-trake} & \pathcode{frontend/src/features/search/Trake*} \\
Video preview / evidence inspection & Fig.~\ref{fig:ui-player} & \pathcode{frontend/src/features/frames} \\
Collaborative workspace / submissions & Fig.~\ref{fig:ui-workspace} & \pathcode{frontend/src/features/workspace}, \pathcode{features/submission} \\
Thunder Compute / S3 execution & Fig.~\ref{fig:architecture}; text below & \pathcode{scripts/build_retrieval_indexes.py}, \pathcode{llm/server} \\
\bottomrule
\end{tabularx}
\end{table}

\begin{figure}[htbp]
\centering
\resizebox{0.98\textwidth}{!}{%
\begin{tikzpicture}[node distance=8mm and 10mm]
  \node[artifact,minimum width=35mm] (s3in) {S3 inputs\\frames, context, transcripts};
  \node[service,minimum width=41mm,right=12mm of s3in] (a6000) {Ephemeral Thunder worker\\one RTX A6000 build profile};
  \node[process,minimum width=34mm,right=12mm of a6000] (validate) {Identity, coverage,\\bundle validation};
  \node[artifact,minimum width=40mm,right=12mm of validate] (s3out) {Versioned S3 publication\\indexes + passed report};
  \draw[flowarrow] (s3in) -- node[above,font=\scriptsize]{selective download} (a6000);
  \draw[flowarrow] (a6000) -- (validate);
  \draw[flowarrow] (validate) -- node[above,font=\scriptsize]{publish only on pass} (s3out);

  \node[service,minimum width=41mm,below=14mm of a6000] (vm) {Optional disposable GPU VM\\Uvicorn inference API :8100};
  \node[process,minimum width=34mm,right=12mm of vm] (edge) {Private Cloudflare\\operator endpoint};
  \node[client,minimum width=40mm,right=12mm of edge] (local) {Local FastAPI + React\\canonical indexes retained};
  \draw[dataarrow] (vm) -- (edge);
  \draw[dataarrow] (edge) -- node[above,font=\scriptsize]{model requests} (local);
  \draw[dataarrow] (s3out.south) |- (local.north);
  \node[note,below=5mm of vm,anchor=north] {Manual lifecycle: create $\rightarrow$ upload $\rightarrow$ health check $\rightarrow$ delete.};
\end{tikzpicture}%
}
\caption{Thunder Compute execution modes. Index construction and optional online inference are separate lifecycles. Canonical serving artifacts are versioned outside the VM, so deleting a GPU instance does not delete retrieval state.}
\label{fig:thunder}
\end{figure}

\paragraph{Release note.}
The runbook's pass-gated publication contract is the intended deployment policy.
The local snapshot documented in Table~\ref{tab:artifacts} does not contain the
promoted build report and has incomplete visual/context provenance fields; those
must be reconciled before an official immutable-bundle claim. The inference
example also demonstrates an L40 while the index-build profile specifies an
A6000. Figure~\ref{fig:thunder} therefore labels GPU choice by workflow rather
than conflating the examples.
